% !TeX program = pdflatex
% =============================================================================
% Arbeitsblatt: Kräfte (Teil 2) (nicht bewertet, Übungsmaterial für den Unterricht)
% Kantonsschule KFR -- Physik, 4. Klasse, 2. FS, Mechanik
%
% Kopiert aus vorlagen/arbeitsblatt-vorlage.tex und um dieselben
% Farb-/Boxdefinitionen ergänzt wie kraefte-teil-1-arbeitsblatt.tex /
% newtons-gesetze-arbeitsblatt.tex. Templates sind selbstständig/nicht DRY,
% siehe CLAUDE.md.
%
% Scope (siehe lernziele.md, Abschnitt "Abgrenzung"): die vier "benannten"
% Kraftarten, die in Kräfte (Teil 1) bewusst ausgeklammert wurden:
% Gewichtskraft (F_G = m*g), Federkraft (Hookesches Gesetz), Normalkraft
% (als Reaktionskraft, Anschluss an das 3. Newtonsche Axiom aus
% newtons-gesetze-arbeitsblatt.tex), Reibungskraft (Haft-/Gleit-/
% Rollreibung, F_R = mu*F_N, inkl. Grenzwinkel mu = tan(alpha)). NICHT Teil
% dieser Einheit: Hydrostatik, Erhaltungssätze, Drehmoment/Statik.
%
% LZ-Nummerierung (LZ1..LZ9) folgt der Dokumentreihenfolge, NICHT der
% Reihenfolge der "Operationalisierten Lernziele" in lernziele.md (dort
% thematisch sortiert). Zuordnung alt -> neu: LZ1->LZ1, LZ2->LZ2, LZ9->LZ3,
% LZ4->LZ4, LZ5->LZ5, LZ3->LZ6, LZ6->LZ7, LZ7->LZ8, LZ8->LZ9. lernziele.md
% wurde entsprechend synchronisiert (Backward-Design-Docx ggf. separat
% nachführen).
%
% Stand: Lernziele, Einleitung, Theorie und Aufgaben fertig.
% =============================================================================
\documentclass[addpoints,11pt,a4paper]{exam}

\usepackage[T1]{fontenc}
\usepackage[utf8]{inputenc}
\usepackage[ngerman]{babel}
\usepackage{lmodern}
\usepackage[a4paper,margin=2.2cm,headheight=14pt]{geometry}
\usepackage{amsmath}
\usepackage{amssymb}
\usepackage{siunitx}
% Hausstil: zusammengesetzte Einheiten immer als echter Bruch (\frac), nicht
% als Schrägstrich inline oder \tfrac. Gilt global für jedes \si/\SI.
\sisetup{per-mode=fraction,fraction-command=\frac}
\usepackage{array,booktabs,tabularx,multirow}
\usepackage{enumitem}
\usepackage{graphicx}
\usepackage{xcolor}
\usepackage{colortbl}
\usepackage{tikz}
\usepackage{physics}
\usepackage{circuitikz}
\usepackage[most]{tcolorbox}
\usepackage{lastpage}

\usetikzlibrary{arrows.meta,calc,angles,quotes}

\setlength{\parindent}{0pt}
\setlength{\parskip}{0.45em}

% -----------------------------
% Farben (Hausfarbschema Physik KFR)
% -----------------------------
\definecolor{titleblue}{HTML}{183A6B}
\definecolor{accentblue}{HTML}{2E6F9E}
\definecolor{lightblue}{HTML}{EAF4FB}
\definecolor{verylightblue}{HTML}{F7FBFE}
\definecolor{lightgray}{HTML}{F4F6F8}
\definecolor{darkgray}{HTML}{4A4A4A}
% Für Diagramme mit zwei verschiedenen Kräften/Linien: titleblue (dunkles
% Blau) + accentorange (Orange) sind weit genug auseinander (auch in
% Graustufen/für Farbenblinde unterscheidbar).
\definecolor{accentorange}{HTML}{D9720C}
% Für Diagramme mit DREI Kräften/Linien: dritte Farbe, die sich sowohl von
% titleblue als auch von accentorange klar unterscheidet.
\definecolor{accentpurple}{HTML}{7B2D8E}
% Theorieblock: eigene Farbe (grün), damit er sich auf einen Blick von den
% blauen Lernziele-/Aufgaben-Boxen abhebt.
\definecolor{theorygreen}{HTML}{2E7D4F}
\definecolor{lighttheorygreen}{HTML}{EAF7EF}
% Gruppenarbeitbox: eigene Farbe (Petrol/Teal), damit Partner-/Gruppenarbeit
% sich auf einen Blick von normalen Einzelaufgaben (blau, taskbox) und vom
% Theorieblock (grün, theoriebox) abhebt.
\definecolor{groupteal}{HTML}{1B7F72}
\definecolor{lightgroupteal}{HTML}{E8F5F3}

% -----------------------------
% Spaltentypen
% -----------------------------
\newcolumntype{Y}{>{\centering\arraybackslash}X}
\newcolumntype{P}[1]{>{\raggedright\arraybackslash}p{#1}}

% -----------------------------
% Boxen
% -----------------------------
\newtcolorbox{titlebox}{
  enhanced,
  colback=titleblue,
  colframe=titleblue,
  boxrule=0pt,
  arc=3mm,
  left=3mm,right=3mm,top=3mm,bottom=3mm
}

% exam.cls rückt die questions-Liste um die Breite von "10." (plus \labelsep)
% ein (Platz für die Nummerierung "1.", "2." ...); wir messen dieselbe
% Distanz hier unabhängig nach.
\newlength{\aufgabenindent}
\settowidth{\aufgabenindent}{10.\hskip\labelsep}

\newtcolorbox{infobox}{
  enhanced,
  breakable,
  colback=lightblue,
  colframe=accentblue,
  boxrule=0.7pt,
  arc=2mm,
  left=5mm,right=5mm,top=2mm,bottom=2mm
}

% theoriebox: wie infobox, aber grün und mit farbigem Titelbalken. Regel:
% hervorgehobene/definierte Begriffe INNERHALB einer theoriebox immer mit
% \textbf{\emph{...}} setzen (fett UND kursiv).
\newtcolorbox{theoriebox}[1][Theorie]{
  enhanced,
  breakable,
  colback=lighttheorygreen,
  colframe=theorygreen,
  title={#1},
  fonttitle=\bfseries,
  boxrule=0.7pt,
  arc=2mm,
  left=5mm,right=5mm,top=2mm,bottom=2mm
}

% taskbox steht direkt in der questions-Liste (vor dem ersten \question),
% also innerhalb von deren Einrückung.
\newtcolorbox{taskbox}[1][]{
  enhanced,
  breakable,
  colback=verylightblue,
  colframe=accentblue,
  title={#1},
  fonttitle=\bfseries,
  boxrule=0.7pt,
  arc=2mm,
  enlarge left by=-\aufgabenindent,
  width=\textwidth,
  left=5mm,right=5mm,top=2mm,bottom=2mm
}

% beispielbox: wie taskbox (blau), aber ausserhalb von questions, deshalb
% KEINE Einrückungskorrektur nötig.
\newtcolorbox{beispielbox}[1][Beispiel]{
  enhanced,
  breakable,
  colback=verylightblue,
  colframe=accentblue,
  title={#1},
  fonttitle=\bfseries,
  boxrule=0.7pt,
  arc=2mm,
  left=5mm,right=5mm,top=2mm,bottom=2mm
}

% Regel für groupbox: JEDE Aufgabe, bei der Schülerinnen/Schüler zu zweit
% oder in einer Gruppe zusammenarbeiten, bekommt diese Box statt der blauen
% taskbox. Strukturell identisch zu taskbox.
\newtcolorbox{groupbox}[1][Gruppenarbeit]{
  enhanced,
  breakable,
  colback=lightgroupteal,
  colframe=groupteal,
  title={#1},
  fonttitle=\bfseries,
  boxrule=0.7pt,
  arc=2mm,
  enlarge left by=-\aufgabenindent,
  width=\textwidth,
  left=5mm,right=5mm,top=2mm,bottom=2mm
}

% hinweisbox: eigene Farbe (Orange), für wichtige Klarstellungen/Warnungen
% vor häufigen Verwechslungen. Steht wie infobox/theoriebox/beispielbox
% ausserhalb von questions, KEINE Einrückungskorrektur nötig.
\definecolor{lightorange}{HTML}{FCEEE0}
\newtcolorbox{hinweisbox}[1][Hinweis]{
  enhanced,
  breakable,
  colback=lightorange,
  colframe=accentorange,
  title={#1},
  fonttitle=\bfseries,
  boxrule=0.7pt,
  arc=2mm,
  left=5mm,right=5mm,top=2mm,bottom=2mm
}

\newtcolorbox{answerbox}{
  breakable,
  colback=white,
  colframe=gray!55,
  boxrule=0.5pt,
  arc=1.5mm,
  left=2mm,right=2mm,top=2mm,bottom=2mm
}

% --- Lösungen ein-/ausblenden -----------------------------------------------
% Schülerexemplar (Standard): \noprintanswers
% Lehrer-/Lösungsexemplar:    \printanswers
\noprintanswers
% \printanswers

% --- Punkte bleiben intern, werden aber nicht gedruckt ----------------------
\pointformat{}

% --- Lösungsbox auf Deutsch, farblich abgesetzt -----------------------------
\renewcommand{\solutiontitle}{\noindent\color{titleblue}\textbf{Lösung:}\enspace}

% --- Kopf-/Fusszeile: Thema/Titel oben, Seitenzahl unten --------------------
\pagestyle{headandfoot}
\cfoot{\textcolor{darkgray}{Seite \thepage\ von \pageref{LastPage}}}

% --- Kopfzeile: Klasse / Datum / Thema / Titel ------------------------------
\newcommand{\Kopfzeile}[4]{%
  \begin{titlebox}
    {\color{white}\sffamily\bfseries\Large #4}
  \end{titlebox}
  \vspace{0.5em}
  \noindent\textbf{Klasse:}~#1 \hfill \textbf{Datum:}~#2 \hfill \textbf{Thema:}~#3
    \hfill \textbf{Lehrperson:}~Loris Keller
  \vspace{0.8em}
  \lhead{\textcolor{darkgray}{#3}}
  \rhead{\textcolor{darkgray}{#4}}
}

% --- Antwortraum: ca. 2.5 cm Platz pro Punkt, als umrandetes Feld -----------
\newlength{\antwortraumlen}
\newcommand{\Antwortraum}[1]{%
  \pgfmathsetlength{\antwortraumlen}{#1*2.5cm}%
  \begin{answerbox}
    \rule{0pt}{\antwortraumlen}%
  \end{answerbox}%
}

% Werte für Kopfzeile zentral setzen
\newcommand{\Klasse}{4a}
\newcommand{\Datum}{TT.MM.JJJJ}
\newcommand{\Thema}{Kräfte (Teil 2)}
\newcommand{\ArbeitsblattTitel}{Arbeitsblatt: Kräfte (Teil 2)}

\begin{document}

\Kopfzeile{\Klasse}{\Datum}{\Thema}{\ArbeitsblattTitel}

% --- Lernziele: aus lernziele.md (Backward-Design, Stufe 1) übernommen,
% aber in Dokumentreihenfolge umnummeriert (siehe Kommentar am Dateianfang).
% ---------------------------------------------------------------------------
\begin{infobox}
\textbf{Lernziele}\\
Am Ende dieser Einheit kann ich \dots
\begin{itemize}[leftmargin=1.2cm,itemsep=0.25em]
  \item die Gewichtskraft $F_G = m\cdot g$ formulieren und Masse und
  Gewichtskraft begrifflich unterscheiden (ortsunabhängig vs.
  ortsabhängig). (LZ1)
  \item die Gewichtskraft eines Körpers an verschiedenen Orten (Erde, Mond)
  berechnen und die Schwerelosigkeit auf der ISS als Zustand des freien
  Falls erklären. (LZ2)
  \item erklären, dass eine Personenwaage keine Masse misst, sondern eine
  Kraft, und diese unter Annahme von $g=\SI{9.81}{\meter\per\second\squared}$
  in eine Masse umrechnet; begründen, wie eine solche Waage für einen
  korrekten Messwert auf dem Mond neu kalibriert werden müsste. (LZ3)
  \item das Hookesche Gesetz $F_H = -k\cdot x$ formulieren und die
  Federkonstante aus einem Kraft-Weg-Diagramm bestimmen. (LZ4)
  \item an einem selbst durchgeführten Federexperiment den linearen
  Zusammenhang zwischen Kraft und Auslenkung überprüfen und grafisch
  auswerten. (LZ5)
  \item die Normalkraft als Reaktionskraft aus dem Kräftegleichgewicht
  senkrecht zur Unterlage bestimmen, auch auf einer geneigten Ebene. (LZ6)
  \item zwischen Haft-, Gleit- und Rollreibung unterscheiden und die
  Reibungskraft mit $F_R = \mu\cdot F_N$ berechnen. (LZ7)
  \item erklären, dass Reibung eine gerichtete, berechenbare Kraft ist und
  nicht ein pauschaler, ungerichteter \glqq Widerstand\grqq{}. (LZ8)
  \item am Grenzwinkel der schiefen Ebene die Beziehung
  $\mu = \tan(\alpha)$ aus dem Kräftegleichgewicht herleiten und mit einem
  phyphox-Experiment (Neigungsmessung) den Haftreibungs- und den
  Gleitreibungskoeffizienten einer Oberfläche experimentell bestimmen.
  (LZ9)
\end{itemize}
\end{infobox}

% --- Einleitung: ISS-Hook, direkt aus lernziele.md übernommen. Bewusst noch
% keine Rechnung im Fliesstext selbst, die folgt strukturiert in der
% Predict-Observe-Explain-Groupbox direkt danach. ---------------------------
\begin{beispielbox}[Einleitung: Schwerelosigkeit auf der ISS]
Die Internationale Raumstation ISS umkreist die Erde in rund
\SI{400}{\kilo\meter} Höhe mit einer Geschwindigkeit von etwa
\SI{28000}{\kilo\meter\per\hour}. Eine einzige Erdumrundung dauert dabei
nur rund \SI{90}{\minute}, die Besatzung erlebt also etwa 16
Sonnenauf- und -untergänge pro Tag.

In Videos von der ISS sieht man die Astronautinnen und Astronauten frei
im Innenraum schweben, zusammen mit Wassertropfen und losen Gegenständen.
Die naheliegende Erklärung lautet oft: \glqq Im Weltall gibt es keine
Schwerkraft.\grqq{} Mit \SI{400}{\kilo\meter} Höhe befindet sich die ISS
aber noch sehr nahe an der Erde, verglichen mit dem Erdradius von rund
\SI{6371}{\kilo\meter}. Ist die Erdanziehungskraft dort wirklich schon
verschwunden?
\end{beispielbox}

\begin{groupbox}[Predict-Observe-Explain: Wirkt auf der ISS noch Schwerkraft?]
\begin{parts}
  \part[2] \textbf{Predict:} Notiert zuerst, ohne zu rechnen, eure eigene
  Vermutung: Wirkt auf Höhe der ISS noch (nahezu) die volle
  Erdanziehungskraft, ist sie deutlich schwächer, oder ist sie praktisch
  null? Begründet kurz.
  \Antwortraum{2}

  \part[3] \textbf{Berechnen:} Die Gravitationskraft nimmt mit dem
  Abstand $r$ vom Erdmittelpunkt gemäss dem quadratischen Abstandsgesetz
  ab:
  \[
    g(r) = g_0\cdot\left(\frac{R_E}{r}\right)^2,
  \]
  wobei $g_0=\SI{9.81}{\meter\per\second\squared}$ der Ortsfaktor an der
  Erdoberfläche und $R_E=\SI{6371}{\kilo\meter}$ der Erdradius ist.
  Berechnet damit gemeinsam $g$ auf Höhe der ISS
  ($r = R_E + \SI{400}{\kilo\meter}$) und gebt das Ergebnis auch als
  Prozentsatz von $g_0$ an.
  \Antwortraum{3}
  \begin{solution}
  \[
    r = R_E + \SI{400}{\kilo\meter} = \SI{6371}{\kilo\meter}+\SI{400}{\kilo\meter} = \SI{6771}{\kilo\meter}.
  \]
  \[
    g(r) = \SI{9.81}{\meter\per\second\squared}\cdot\left(\frac{\SI{6371}{\kilo\meter}}{\SI{6771}{\kilo\meter}}\right)^2
    \approx \SI{9.81}{\meter\per\second\squared}\cdot(0.9409)^2
  \]
  \[
    \approx \SI{9.81}{\meter\per\second\squared}\cdot0.8853 \approx \SI{8.68}{\meter\per\second\squared}.
  \]
  Als Prozentsatz von $g_0$: $\dfrac{\SI{8.68}{\meter\per\second\squared}}{\SI{9.81}{\meter\per\second\squared}} \approx 0.885 \; \widehat{=}\; \SI{88.5}{\percent}$.
  \end{solution}

  \part[3] \textbf{Explain:} Auf der ISS wirkt also immer noch rund
  \SI{88.5}{\percent} der Erdanziehungskraft von der Erdoberfläche, die
  Astronautin schwebt also nicht, weil \glqq keine Schwerkraft
  herrscht\grqq{}. Erkläre stattdessen, warum sie trotzdem schwebt. Denke
  dabei an den waagrechten Wurf: Was passiert mit einem geworfenen
  Gegenstand, wenn seine Wurfweite so gross wird, dass die gekrümmte
  Erdoberfläche unter ihm \glqq wegfällt\grqq{}, bevor er landet?
  \Antwortraum{3}
  \begin{solution}
  Die ISS bewegt sich so schnell seitwärts, dass sie während des
  Fallens ständig \glqq an der gekrümmten Erdoberfläche vorbeifällt\grqq{},
  genau wie beim waagrechten Wurf, nur mit so hoher Geschwindigkeit, dass
  die Bahnkurve nie am Boden ankommt, sondern sich zu einem Kreis um die
  Erde schliesst. Die ISS befindet sich also im permanenten freien Fall.
  Da sich Astronautin und Raumstation gemeinsam mit derselben Beschleunigung
  im freien Fall befinden, übt die Station keine Normalkraft mehr auf die
  Astronautin aus, es entsteht das Gefühl der Schwerelosigkeit, obwohl die
  Gravitationskraft selbst kaum reduziert ist.
  \end{solution}
\end{parts}
\end{groupbox}

\begin{hinweisbox}[Praktisch: Mini-Experiment mit dem Federkraftmesser]
Bevor die Gewichtskraft formal eingeführt wird, ein kurzer
Demonstrationsversuch: Ein bekanntes Gewichtsstück wird an einen
Federkraftmesser gehängt (Messprinzip siehe Kräfte Teil 1). Der
abgelesene Kraftwert in Newton ist genau die Gewichtskraft dieses
Gewichtsstücks, die direkte Überleitung zur folgenden Theoriebox.
\end{hinweisbox}

\begin{theoriebox}[Die Gewichtskraft $F_G = m\cdot g$]
Jeder Körper mit Masse $m$ wird von der Erde angezogen. Diese Anziehung
ist die \textbf{\emph{Gewichtskraft}} $F_G$:
\[
  \boxed{F_G = m\cdot g}, \qquad \text{Einheit von } F_G\colon\ \si{\newton}.
\]
Dabei ist $g$ der \textbf{\emph{Ortsfaktor}} (auch Fallbeschleunigung
genannt), auf der Erdoberfläche $g_0\approx\SI{9.81}{\meter\per\second\squared}$.
Die Gewichtskraft zeigt stets senkrecht nach unten, das heisst in
Richtung des Erdmittelpunkts.

\textbf{\emph{Masse}} $m$ (Einheit $\si{\kilo\gram}$) und
\textbf{\emph{Gewichtskraft}} $F_G$ (Einheit $\si{\newton}$) sind nicht
dasselbe:
\begin{itemize}[leftmargin=1.2cm,itemsep=0.25em]
  \item Die \textbf{\emph{Masse}} ist ortsunabhängig: Ein Körper mit
  $\SI{60}{\kilo\gram}$ hat auf der Erde, auf dem Mond und im leeren
  Weltraum immer dieselbe Masse.
  \item Die \textbf{\emph{Gewichtskraft}} ist ortsabhängig: Sie hängt über
  $g$ vom Himmelskörper (oder, wie bei der ISS, von der Höhe über der
  Erdoberfläche) ab, an dem sich der Körper gerade befindet.
\end{itemize}
\end{theoriebox}

\begin{hinweisbox}[Hinweis: \glqq wiegen\grqq{} ist nicht gleich Masse]
Im Alltag sagt man oft \glqq er wiegt \SI{60}{\kilo\gram}\grqq{}, dabei
wird eigentlich die Masse gemeint, nicht die Gewichtskraft. Physikalisch
korrekt wiegt (also: hat als Gewichtskraft) eine Person mit
$m=\SI{60}{\kilo\gram}$ auf der Erde $F_G\approx\SI{588.6}{\newton}$, auf
dem Mond aber nur $F_G\approx\SI{97.2}{\newton}$, obwohl ihre Masse an
beiden Orten unverändert $\SI{60}{\kilo\gram}$ bleibt. Ein Astronaut wiegt
auf dem Mond also weniger, hat aber nicht weniger Masse.
\end{hinweisbox}

\begin{theoriebox}[Die Gewichtskraft an anderen Orten]
Der Ortsfaktor $g$ ist nicht überall gleich gross, sondern hängt von der
Masse und dem Radius des jeweiligen Himmelskörpers ab:

\begin{center}
\begin{tabular}{@{}lll@{}}
\toprule
\textbf{Ort} & \textbf{Ortsfaktor $g$} & \textbf{$F_G$ für $m=\SI{60}{\kilo\gram}$} \\
\midrule
Erdoberfläche & $\SI{9.81}{\meter\per\second\squared}$ & $\SI{588.6}{\newton}$ \\
Höhe der ISS ($\approx\SI{400}{\kilo\meter}$) & $\approx\SI{8.68}{\meter\per\second\squared}$ & $\approx\SI{520.8}{\newton}$ \\
Mond & $\SI{1.62}{\meter\per\second\squared}$ & $\SI{97.2}{\newton}$ \\
Mars & $\SI{3.71}{\meter\per\second\squared}$ & $\SI{222.6}{\newton}$ \\
\bottomrule
\end{tabular}
\end{center}

Auf der ISS ist $g$ zwar gegenüber der Erdoberfläche reduziert (siehe
Predict-Observe-Explain oben), aber bei Weitem nicht so stark wie auf dem
Mond. Der Grund für die Schwerelosigkeit auf der ISS ist deshalb nicht ein
kleines $g$, sondern der permanente freie Fall.
\end{theoriebox}

\begin{questions}
\begin{taskbox}[Übung: Gewichtskraft berechnen]
\question[4] Eine Person hat eine Masse von $m=\SI{72}{\kilo\gram}$.

\begin{parts}
  \part[2] Berechne ihre Gewichtskraft auf der Erde.
  \Antwortraum{2}
  \begin{solution}
  \[
    F_G = m\cdot g_0 = \SI{72}{\kilo\gram}\cdot\SI{9.81}{\meter\per\second\squared} \approx \SI{706.3}{\newton}.
  \]
  \end{solution}

  \part[2] Berechne ihre Gewichtskraft auf dem Mond.
  \Antwortraum{2}
  \begin{solution}
  \[
    F_G = m\cdot g_{\text{Mond}} = \SI{72}{\kilo\gram}\cdot\SI{1.62}{\meter\per\second\squared} \approx \SI{116.6}{\newton}.
  \]
  \end{solution}
\end{parts}

\question[3] \glqq Ein Astronaut wiegt auf dem Mars weniger als auf der
Erde. Hat er dort nur weniger Masse, hat er nur eine kleinere
Gewichtskraft, oder beides?\grqq{}

\begin{parts}
  \part[3] Beantworte die Frage in ein bis zwei Sätzen und begründe sie
  mit den Begriffen Masse und Gewichtskraft. Nutze zur Begründung die
  Werte für $g$ aus der Tabelle in der vorherigen Theoriebox.
  \Antwortraum{3}
  \begin{solution}
  Nur die Gewichtskraft ist kleiner, nicht die Masse. Die Masse des
  Astronauten bleibt auf dem Mars exakt gleich (sie ist ortsunabhängig),
  aber der Mars hat mit $g_{\text{Mars}}=\SI{3.71}{\meter\per\second\squared}$
  einen kleineren Ortsfaktor als die Erde
  ($g_0=\SI{9.81}{\meter\per\second\squared}$), deshalb ist
  $F_G=m\cdot g$ dort kleiner, genau wie beim Mond-Beispiel aus dem
  Hinweis oben, nur mit einem anderen Zahlenwert.
  \end{solution}
\end{parts}

\question[4] Dieselbe Person aus der ersten Teilaufgabe ($m=\SI{72}{\kilo\gram}$)
befindet sich nun an Bord der ISS.

\begin{parts}
  \part[2] Berechne ihre Gewichtskraft dort, mit
  $g_{\text{ISS}}\approx\SI{8.68}{\meter\per\second\squared}$ aus der
  Predict-Observe-Explain-Aufgabe.
  \Antwortraum{2}
  \begin{solution}
  \[
    F_G = m\cdot g_{\text{ISS}} = \SI{72}{\kilo\gram}\cdot\SI{8.68}{\meter\per\second\squared} \approx \SI{624.96}{\newton}.
  \]
  \end{solution}

  \part[2] Bedeutet das Ergebnis aus a), dass auf die Person an Bord der
  ISS praktisch keine Kraft mehr wirkt? Begründe kurz.
  \Antwortraum{2}
  \begin{solution}
  Nein: Die Gewichtskraft ist mit rund $\SI{625}{\newton}$ fast so gross
  wie auf der Erde ($\SI{706.3}{\newton}$). Die Person schwebt trotzdem,
  weil sie sich zusammen mit der Station im freien Fall befindet, nicht
  weil keine Kraft wirkt.
  \end{solution}
\end{parts}
\end{taskbox}
\end{questions}

% --- Hausaufgabe (Rechercheauftrag), wie in lernziele.md vorgeschlagen
% ("Hausaufgabe nach L1"), deshalb hier direkt nach der Gewichtskraft
% platziert, vor der Federkraft. LZ3 (alt LZ9). ------------------------------
\begin{questions}
\begin{taskbox}[Hausaufgabe: Wie funktioniert eine Personenwaage?]
\question[8] Eine haushaltsübliche Personenwaage zeigt dir eine Zahl in
Kilogramm an, wenn du daraufstehst. Recherchiere (Internet oder
KI-Chatbot, mit kritischer Prüfung der Quelle) und beantworte:

\begin{parts}
  \part[3] Wie misst eine moderne Personenwaage überhaupt: über eine
  mechanische Feder oder über einen elektronischen
  Dehnungsmessstreifensensor? Beschreibe das Messprinzip kurz.
  \Antwortraum{3}
  \begin{solution}
  Moderne Haushaltswaagen verwenden meist Dehnungsmessstreifensensoren
  (auch Wägezellen genannt): Das Gewicht verformt ein Metallelement
  minimal, diese Verformung ändert den elektrischen Widerstand des
  aufgeklebten Dehnungsmessstreifens, und aus dieser Widerstandsänderung
  berechnet die Elektronik einen Kraftwert. Ältere/mechanische Waagen
  nutzen stattdessen eine Feder, ähnlich dem Federkraftmesser.
  \end{solution}

  \part[2] Was misst der Sensor der Waage physikalisch tatsächlich, und
  wie rechnet die Waage daraus die angezeigte Masse in Kilogramm aus?
  \Antwortraum{2}
  \begin{solution}
  Der Sensor misst eine Kraft (die Gewichtskraft $F_G$), nicht direkt
  eine Masse. Die Elektronik rechnet intern mit $m = F_G / g$ um, wobei
  $g=\SI{9.81}{\meter\per\second\squared}$ (der Erdoberflächenwert) fest
  einprogrammiert ist, und zeigt das Ergebnis als Masse in Kilogramm an.
  \end{solution}

  \part[3] Wie müsste eine solche Waage verändert (neu kalibriert)
  werden, damit sie auf dem Mond weiterhin die korrekte Masse anzeigt?
  \Antwortraum{3}
  \begin{solution}
  Der fest einprogrammierte Wert $g=\SI{9.81}{\meter\per\second\squared}$
  müsste durch $g_{\text{Mond}}=\SI{1.62}{\meter\per\second\squared}$
  ersetzt werden, damit die Waage weiterhin korrekt $m=F_G/g$ berechnet.
  Ohne diese Anpassung würde die unveränderte Waage auf dem Mond eine viel
  zu kleine Masse anzeigen, da dieselbe Masse dort nur eine viel kleinere
  Gewichtskraft erzeugt, aus der die Waage (mit dem falschen, zu grossen
  $g_0$) fälschlich auf eine kleinere Masse schliessen würde.
  \end{solution}
\end{parts}
\end{taskbox}
\end{questions}

\begin{theoriebox}[Die Federkraft: das Hookesche Gesetz]
Wie in Kräfte (Teil 1) beobachtet, dehnt sich eine Feder umso weiter, je
stärker eine Kraft an ihr zieht, und diese Dehnung ist proportional zur
Kraft. Dieser Zusammenhang heisst \textbf{\emph{Hookesches Gesetz}}:
\[
  \boxed{F_H = -k\cdot x}, \qquad \text{Einheit der \textbf{\emph{Federkonstante}} } k\colon\ \si{\newton\per\meter}.
\]
Dabei ist $x$ die Auslenkung (Dehnung oder Stauchung, Einheit
$\si{\meter}$) aus der \textbf{\emph{Ruhelage}} der Feder, und $k$ die
\textbf{\emph{Federkonstante}}, ein Mass dafür, wie \glqq steif\grqq{}
eine Feder ist (grosses $k$: harte Feder, kleines $k$: weiche Feder). Das
Minuszeichen bedeutet: $F_H$ ist eine \textbf{\emph{rückstellende
Kraft}}, sie zeigt immer zurück in Richtung der Ruhelage, egal ob die
Feder gedehnt oder gestaucht ist. Genau deshalb federt eine ausgelenkte
Feder immer wieder zur selben Ruhelage zurück: Je weiter sie ausgelenkt
wird, desto stärker zieht (oder drückt) sie zurück.

Trägt man den Betrag der Kraft $F$ über der Auslenkung $x$ auf, ergibt
sich eine Gerade durch den Ursprung, deren Steigung genau $k$ ist:

\begin{center}
\begin{tikzpicture}[scale=1.4]
  \draw[->,thick] (0,0) -- (5.5,0) node[right] {\scriptsize $x$};
  \draw[->,thick] (0,0) -- (0,4.5) node[above] {\scriptsize $F$};
  \draw[thick,titleblue] (0,0) -- (5,4);
  \draw[dashed,gray] (3,0) -- (3,2.4) -- (0,2.4);
  \node[below] at (3,0) {\scriptsize $x_1$};
  \node[left] at (0,2.4) {\scriptsize $F_1$};
\end{tikzpicture}\\[0.2em]
{\scriptsize Kraft-Weg-Diagramm einer Feder: Die Steigung der Geraden ist
die Federkonstante $k=\dfrac{F_1}{x_1}$.}
\end{center}
\end{theoriebox}

\begin{beispielbox}[Beispiel: Federkonstante aus zwei Messwerten]
Eine Feder dehnt sich unter einer Kraft von $F=\SI{8}{\newton}$ um
$x=\SI{4}{\centi\meter}$.

\textbf{Federkonstante:}
\[
  k = \frac{F}{x} = \frac{\SI{8}{\newton}}{\SI{0.04}{\meter}} = \SI{200}{\newton\per\meter}.
\]
\textbf{Vorhersage für eine neue Kraft} $F_2=\SI{15}{\newton}$:
\[
  x_2 = \frac{F_2}{k} = \frac{\SI{15}{\newton}}{\SI{200}{\newton\per\meter}} = \SI{0.075}{\meter} = \SI{7.5}{\centi\meter}.
\]
\end{beispielbox}

\begin{questions}
\begin{taskbox}[Übung: mit der Federkonstante rechnen]
\question[6] Eine andere Feder dehnt sich unter einer Kraft von
$F=\SI{9}{\newton}$ um $x=\SI{3}{\centi\meter}$.

\begin{parts}
  \part[3] Berechne die Federkonstante $k$ dieser Feder, und gib die
  Einheit korrekt an.
  \Antwortraum{3}
  \begin{solution}
  \[
    k = \frac{F}{x} = \frac{\SI{9}{\newton}}{\SI{0.03}{\meter}} = \SI{300}{\newton\per\meter}.
  \]
  \end{solution}

  \part[3] Berechne, um wie viel sich dieselbe Feder unter einer Kraft
  von $F=\SI{21}{\newton}$ dehnt.
  \Antwortraum{3}
  \begin{solution}
  \[
    x = \frac{F}{k} = \frac{\SI{21}{\newton}}{\SI{300}{\newton\per\meter}} = \SI{0.07}{\meter} = \SI{7}{\centi\meter}.
  \]
  \end{solution}
\end{parts}

\question[3] Für eine dritte Feder wurde das folgende Kraft-Weg-Diagramm
gemessen.

\begin{center}
\begin{tikzpicture}[scale=0.9]
  \draw[gray!35,thin] (0,0) grid (6,6);
  \draw[->,thick] (0,0) -- (6.4,0) node[right] {\scriptsize $x$ in cm};
  \draw[->,thick] (0,0) -- (0,6.4) node[above] {\scriptsize $F$ in N};
  \foreach \x/\lbl in {1/1,2/2,3/3,4/4,5/5,6/6}
    \node[below] at (\x,0) {\scriptsize \lbl};
  \foreach \y/\lbl in {1/2,2/4,3/6,4/8,5/10,6/12}
    \node[left] at (0,\y) {\scriptsize \lbl};
  \draw[thick,titleblue] (0,0) -- (5,5);
  \draw[dashed,gray] (5,0) -- (5,5) -- (0,5);
\end{tikzpicture}
\end{center}

\begin{parts}
  \part[3] Lies ein Wertepaar $(x,F)$ auf der Geraden ab und bestimme
  damit die Federkonstante $k$ dieser Feder.
  \Antwortraum{3}
  \begin{solution}
  Abgelesen z.\,B. bei $x=\SI{5}{\centi\meter}=\SI{0.05}{\meter}$,
  $F=\SI{10}{\newton}$:
  \[
    k = \frac{F}{x} = \frac{\SI{10}{\newton}}{\SI{0.05}{\meter}} = \SI{200}{\newton\per\meter}.
  \]
  \end{solution}
\end{parts}
\end{taskbox}
\end{questions}

% --- Hands-On-Experiment (M5), wie in lernziele.md vorgeschlagen (L2),
% ergänzend PhET-Simulation "Hooke's Law". LZ5. ------------------------------
\begin{groupbox}[Federexperiment: die Federkonstante bestimmen]
Bearbeitet die folgende Aufgabe zu zweit oder in einer Kleingruppe mit
einer echten Feder und einem Satz Gewichtsstücke bekannter Kraft (z.\,B.
$\SI{1}{\newton}$, $\SI{2}{\newton}$, $\SI{3}{\newton}$, $\SI{4}{\newton}$).
Ergänzend könnt ihr eure Ergebnisse mit der PhET-Simulation
\glqq Hooke's Law\grqq{} vergleichen.

\begin{parts}
  \part[3] \textbf{Messen:} Hängt der Reihe nach die Gewichtsstücke an
  eure Feder und messt jeweils die Auslenkung $x$ mit einem Lineal.
  Tragt eure Messwerte in die Tabelle ein.
  \begin{center}
  \begin{tabular}{@{}lp{2.8cm}p{2.8cm}@{}}
  \toprule
  \textbf{Messung} & \textbf{Kraft $F$} & \textbf{Auslenkung $x$} \\
  \midrule
  1 & & \\[0.6cm]
  2 & & \\[0.6cm]
  3 & & \\[0.6cm]
  4 & & \\[0.6cm]
  \bottomrule
  \end{tabular}
  \end{center}

  \part[3] \textbf{Auswerten:} Zeichnet euer Kraft-Weg-Diagramm ($F$ über
  $x$) auf dem folgenden Gitter.
  \begin{center}
  \begin{tikzpicture}[scale=1.3]
    \draw[step=1cm,gray!35,thin] (0,0) grid (7,7);
    \draw[->,thick] (0,0) -- (7.4,0) node[right] {\scriptsize $x$};
    \draw[->,thick] (0,0) -- (0,7.4) node[above] {\scriptsize $F$};
  \end{tikzpicture}
  \end{center}

  \part[4] \textbf{Federkonstante:} Bestimmt aus der Steigung eurer
  Geraden ($k=\Delta F/\Delta x$) die Federkonstante eurer Feder, und
  vergleicht euer Ergebnis mit der PhET-Simulation \glqq Hooke's
  Law\grqq{}: Bestätigt die Simulation den linearen Zusammenhang aus
  eurer eigenen Messung?
  \Antwortraum{4}
  \begin{solution}
  Individuelle Messwerte, von der Lehrperson auf einen linearen
  Zusammenhang zu kontrollieren. Beispielhafte Auswertung für eine Feder
  mit $k=\SI{20}{\newton\per\meter}$: $F=\SI{1}{\newton}\to x=\SI{5}{\centi\meter}$,
  $F=\SI{2}{\newton}\to x=\SI{10}{\centi\meter}$,
  $F=\SI{3}{\newton}\to x=\SI{15}{\centi\meter}$,
  $F=\SI{4}{\newton}\to x=\SI{20}{\centi\meter}$, alle Punkte liegen auf
  einer Geraden durch den Ursprung mit Steigung
  $k=\dfrac{\SI{4}{\newton}}{\SI{0.2}{\meter}}=\SI{20}{\newton\per\meter}$.
  Die PhET-Simulation sollte denselben linearen Zusammenhang zeigen (die
  angezeigte Feder gehorcht demselben Hookeschen Gesetz).
  \end{solution}
\end{parts}
\end{groupbox}

% --- Clicker-ConcepTest (M2) zur Aktiv-Passiv-Lernschwierigkeit, als
% informelle Erkundung vor der Normalkraft-Theorie, wie in lernziele.md
% vorgeschlagen (informell mit PhET "Friction" als Einstieg gedacht, hier
% direkt als Konzepttest zur Normalkraft/Aktiv-Passiv-Schema). --------------
\begin{questions}
\begin{taskbox}[Clicker-ConcepTest: Übt der Tisch eine Kraft aus?]
\question[4] Ein Buch liegt ruhig auf einem Tisch. Die Gewichtskraft des
Buchs zieht es nach unten, trotzdem bewegt es sich nicht.

\begin{parts}
  \part[2] Kreuze an: Übt der Tisch dabei aktiv eine Kraft auf das Buch
  aus, oder \glqq leistet\grqq{} er nur passiv Widerstand, ohne selbst
  eine Kraft auszuüben? Begründe kurz.
  \begin{itemize}[leftmargin=1.2cm,itemsep=0.15em]
    \item[$\square$] A: Der Tisch übt aktiv eine Kraft nach oben aus, die
    die Gewichtskraft des Buchs genau ausgleicht.
    \item[$\square$] B: Der Tisch übt keine eigene Kraft aus, er
    \glqq blockiert\grqq{} das Buch nur passiv.
  \end{itemize}
  \Antwortraum{2}
  \begin{solution}
  Richtig ist \textbf{A}. Ein starrer Körper wie ein Tisch verformt sich
  unter der Last minimal (auch wenn das Auge es nicht sieht) und übt
  dadurch tatsächlich eine aktive, gerichtete Gegenkraft aus, die
  \textbf{\emph{Normalkraft}} $F_N$. Ohne diese aktive Gegenkraft würde
  das Buch dem 2. Newtonschen Gesetz folgend beschleunigt nach unten
  fallen.
  \end{solution}

  \part[2] Was würde passieren, wenn der Tisch (z.\,B. weil er zu
  schwach gebaut ist) diese Kraft nicht vollständig aufbringen könnte?
  \Antwortraum{2}
  \begin{solution}
  Dann wäre die Gewichtskraft grösser als die Normalkraft des Tisches, es
  läge kein Kräftegleichgewicht mehr vor, und das Buch (bzw. der Tisch
  selbst) würde sich nach unten in Bewegung setzen, im Extremfall bricht
  der Tisch durch.
  \end{solution}
\end{parts}
\end{taskbox}
\end{questions}

\begin{theoriebox}[Die Normalkraft]
Nach dem 3. Newtonschen Gesetz (Wechselwirkungsgesetz, siehe Newtons
Gesetze) übt jede Unterlage, auf der ein Körper aufliegt, eine
Reaktionskraft senkrecht zu ihrer eigenen Oberfläche aus: die
\textbf{\emph{Normalkraft}} $F_N$ (\glqq normal\grqq{} bedeutet hier
\glqq senkrecht\grqq{}).

Wichtig: $F_N$ ist \textbf{nicht} automatisch gleich der Gewichtskraft.
Ihr tatsächlicher Betrag ergibt sich immer aus dem
\textbf{\emph{Kräftegleichgewicht senkrecht zur Unterlage}}, das heisst,
man muss alle Kräfte in dieser Richtung betrachten:
\begin{itemize}[leftmargin=1.2cm,itemsep=0.25em]
  \item Auf einer \textbf{\emph{waagrechten}} Unterlage ohne weitere
  vertikale Kräfte gilt $F_N = F_G$.
  \item Wirkt zusätzlich eine weitere vertikale Kraft (z.\,B. jemand
  drückt von oben), muss diese im Gleichgewicht mitgezählt werden.
  \item Auf einer \textbf{\emph{geneigten Ebene}} (siehe Kräfte Teil 1)
  ist $F_N$ gleich der Normalkomponente der Gewichtskraft:
  \[
    \boxed{F_N = F_G\cdot\cos(\alpha)}.
  \]
  In Kräfte (Teil 1) wurde diese Komponente bewusst
  \glqq Normalkomponente\grqq{} genannt, weil das 3. Newtonsche Gesetz
  noch nicht bekannt war. Jetzt wisst ihr: Diese Komponente ist genau die
  Reaktionskraft der Ebene, also die Normalkraft $F_N$ selbst.
\end{itemize}
\end{theoriebox}

\begin{beispielbox}[Beispiel: Normalkraft mit zusätzlicher Kraft]
Eine Kiste mit Gewichtskraft $F_G=\SI{500}{\newton}$ steht auf einem
waagrechten Boden. Eine Person drückt zusätzlich mit $F_P=\SI{150}{\newton}$
senkrecht von oben auf die Kiste.

\textbf{Kräftegleichgewicht senkrecht (die Kiste bewegt sich nicht):}
\[
  F_N - F_G - F_P = \SI{0}{\newton}.
\]
\textbf{Nach $F_N$ auflösen:}
\[
  \boxed{F_N = F_G + F_P} = \SI{500}{\newton}+\SI{150}{\newton} = \SI{650}{\newton}.
\]
Die Normalkraft ist hier also grösser als die Gewichtskraft allein.
\end{beispielbox}

\begin{questions}
\begin{taskbox}[Übung: Normalkraft bestimmen]
\question[8] Zwei Situationen mit unterschiedlichen Normalkräften.

\begin{parts}
  \part[3] Eine Kiste mit $F_G=\SI{500}{\newton}$ steht auf waagrechtem
  Boden. Eine Person drückt zusätzlich mit $F_P=\SI{150}{\newton}$ von
  oben darauf (wie im Beispiel oben, aber jetzt selbst rechnen).
  Bestimme $F_N$ über das Kräftegleichgewicht.
  \Antwortraum{3}
  \begin{solution}
  \[
    F_N = F_G + F_P = \SI{500}{\newton}+\SI{150}{\newton} = \SI{650}{\newton}.
  \]
  \end{solution}

  \part[3] Der Snowboarder aus Kräfte (Teil 1), $F_G=\SI{650}{\newton}$,
  steht auf einer Piste mit Neigungswinkel $\alpha=\ang{20}$. Bestimme
  seine Normalkraft $F_N$ (dort noch \glqq Normalkomponente\grqq{}
  genannt).
  \Antwortraum{3}
  \begin{solution}
  \[
    F_N = F_G\cdot\cos(\alpha) = \SI{650}{\newton}\cdot\cos(\ang{20}) \approx \SI{650}{\newton}\cdot0.940 \approx \SI{610.8}{\newton}.
  \]
  \end{solution}

  \part[2] Ist die Normalkraft also immer gleich der Gewichtskraft?
  Begründe mit Verweis auf die Teilaufgaben a) und b).
  \Antwortraum{2}
  \begin{solution}
  Nein. In a) ist $F_N=\SI{650}{\newton}$ grösser als $F_G=\SI{500}{\newton}$
  (zusätzliche Druckkraft), in b) ist $F_N\approx\SI{610.8}{\newton}$
  kleiner als $F_G=\SI{650}{\newton}$ (geneigte Ebene). $F_N$ ergibt sich
  immer erst aus dem Kräftegleichgewicht senkrecht zur Unterlage, nicht
  automatisch aus $F_G$ allein.
  \end{solution}
\end{parts}
\end{taskbox}
\end{questions}

\begin{hinweisbox}[Praktisch: Erste Erkundung mit PhET \glqq Friction\grqq{}]
Bevor die Reibung formal eingeführt wird, ein kurzer informeller
Einstieg: In der PhET-Simulation \glqq Friction\grqq{} lässt sich ein
Block über verschiedene Oberflächen schieben, während die Simulation die
Temperatur an der Kontaktfläche und die Bewegung sichtbar macht.
Experimentiert kurz selbst damit: Bei welcher Oberfläche braucht es mehr
Kraft, den Block überhaupt erst in Bewegung zu setzen?
\end{hinweisbox}

\begin{theoriebox}[Reibung: Haft-, Gleit- und Rollreibung]
Reibung entsteht, wenn sich zwei Oberflächen berühren und sich relativ
zueinander bewegen wollen oder bewegen. Man unterscheidet drei Arten:

\begin{itemize}[leftmargin=1.2cm,itemsep=0.3em]
  \item \textbf{\emph{Haftreibung}}: verhindert, dass ein ruhender Körper
  überhaupt erst anfängt zu rutschen (z.\,B. ein Bücherstapel, den man
  erst mit genügend Kraft in Bewegung setzen muss).
  \item \textbf{\emph{Gleitreibung}}: wirkt, sobald sich der Körper
  bereits über die Oberfläche bewegt (gleitet/rutscht).
  \item \textbf{\emph{Rollreibung}}: wirkt bei rollenden Körpern (Rad,
  Kugel) und ist deutlich kleiner als Gleitreibung, deshalb rollen Räder
  so viel leichter, als dass Dinge rutschen.
\end{itemize}

In allen drei Fällen berechnet sich die Reibungskraft aus der Normalkraft
$F_N$ (nicht aus der Gewichtskraft direkt, siehe vorherige Theoriebox):
\[
  \boxed{F_R = \mu\cdot F_N},
\]
wobei $\mu$ der (einheitenlose) \textbf{\emph{Reibungskoeffizient}} ist,
mit eigenem Wert für Haft- ($\mu_H$), Gleit- ($\mu_G$) und Rollreibung
($\mu_R$). Für dieselbe Materialpaarung gilt praktisch immer
$\mu_H > \mu_G > \mu_R$, deshalb ist es schwerer, einen Gegenstand erst
in Bewegung zu setzen, als ihn danach in Bewegung zu halten.

\begin{center}
\begin{tabular}{@{}lll@{}}
\toprule
\textbf{Materialpaarung} & \textbf{$\mu_H$} & \textbf{$\mu_G$} \\
\midrule
Holz auf Holz & $0.5$ & $0.3$ \\
Gummi auf trockenem Asphalt & $0.8$ & $0.7$ \\
Stahl auf Stahl & $0.15$ & $0.12$ \\
Schlittschuh auf Eis & $0.03$ & $0.01$ \\
\bottomrule
\end{tabular}
\end{center}

Zum Vergleich: Ein rollendes Autorad auf Asphalt hat eine Rollreibung von
nur $\mu_R\approx 0.01$ bis $0.02$, deutlich kleiner als die Gleitreibung
von Gummi auf Asphalt. Das ist einer der Gründe, weshalb Fahrzeuge Räder
statt Kufen verwenden.
\end{theoriebox}

\begin{hinweisbox}[Hinweis: Reibung ist eine gerichtete Kraft, kein pauschaler \glqq Widerstand\grqq{}]
Im Alltag wird Reibung oft als vager, ungerichteter \glqq Widerstand\grqq{}
verstanden, so als würde sie einfach \glqq bremsen\grqq{}, egal in welche
Richtung. Physikalisch ist die Reibungskraft $F_R$ aber eine ganz normale
Kraft mit Betrag \emph{und} Richtung: Sie wirkt immer genau entgegen der
Richtung, in die sich der Körper relativ zur Unterlage bewegt (oder zu
bewegen versucht), und ihr Betrag lässt sich mit $F_R=\mu\cdot F_N$ exakt
berechnen, nicht nur qualitativ abschätzen.
\end{hinweisbox}

\begin{questions}
\begin{taskbox}[Übung: Reibungskraft berechnen]
\question[9] Eine Kiste mit Masse $m=\SI{50}{\kilo\gram}$ steht auf
einem waagrechten Holzboden ($\mu_H=0.5$, $\mu_G=0.3$, siehe Tabelle
oben).

\begin{parts}
  \part[3] Berechne zuerst die Normalkraft, danach die maximale
  Haftreibungskraft, die überwunden werden muss, damit die Kiste
  überhaupt zu rutschen beginnt.
  \Antwortraum{3}
  \begin{solution}
  \[
    F_N = F_G = m\cdot g = \SI{50}{\kilo\gram}\cdot\SI{9.81}{\meter\per\second\squared} \approx \SI{490.5}{\newton}.
  \]
  \[
    F_{R,H,\max} = \mu_H\cdot F_N = 0.5\cdot\SI{490.5}{\newton} \approx \SI{245.3}{\newton}.
  \]
  \end{solution}

  \part[3] Berechne die Gleitreibungskraft, sobald die Kiste bereits
  rutscht.
  \Antwortraum{3}
  \begin{solution}
  \[
    F_{R,G} = \mu_G\cdot F_N = 0.3\cdot\SI{490.5}{\newton} \approx \SI{147.15}{\newton}.
  \]
  \end{solution}

  \part[3] Vergleiche die beiden Ergebnisse: Warum ist es schwerer, die
  Kiste in Bewegung zu \textbf{setzen}, als sie danach in Bewegung zu
  \textbf{halten}?
  \Antwortraum{3}
  \begin{solution}
  Die maximale Haftreibungskraft ($\approx\SI{245.3}{\newton}$) ist
  grösser als die Gleitreibungskraft ($\approx\SI{147.15}{\newton}$),
  weil $\mu_H>\mu_G$ gilt. Man muss also zuerst die grössere
  Haftreibungskraft überwinden, um die Kiste überhaupt in Bewegung zu
  setzen; ist sie erst einmal in Bewegung, reicht die kleinere
  Gleitreibungskraft, um sie in Bewegung zu halten.
  \end{solution}
\end{parts}

\question[3] Ordne jeder Situation die passende Reibungsart zu:
\textbf{(A)} Haftreibung, \textbf{(B)} Gleitreibung, \textbf{(C)} Rollreibung.

\begin{parts}
  \part[1] Ein Fahrradreifen rollt über die Strasse.
  \Antwortraum{1}
  \begin{solution}
  \textbf{C} (Rollreibung, deutlich kleiner als Gleitreibung, deshalb
  rollen Räder so viel leichter, als dass Dinge rutschen).
  \end{solution}

  \part[1] Ein Schlitten gleitet bereits über den Schnee.
  \Antwortraum{1}
  \begin{solution}
  \textbf{B} (Gleitreibung, der Schlitten bewegt sich bereits über die
  Unterlage).
  \end{solution}

  \part[1] Ein Schrank steht ruhig auf dem Boden und soll gerade erst
  verschoben werden.
  \Antwortraum{1}
  \begin{solution}
  \textbf{A} (Haftreibung, sie verhindert, dass der Schrank überhaupt
  erst zu rutschen beginnt).
  \end{solution}
\end{parts}
\end{taskbox}
\end{questions}

\begin{theoriebox}[Der Grenzwinkel der schiefen Ebene: $\mu_H = \tan(\alpha)$]
Legt man einen Körper auf eine schiefe Ebene und vergrössert langsam den
Neigungswinkel $\alpha$, bleibt er zunächst liegen (Haftreibung hält ihn
fest), bis er ab einem bestimmten \textbf{\emph{Grenzwinkel}} zu rutschen
beginnt. Genau in diesem Moment ist die maximale Haftreibungskraft gerade
so gross wie die Hangabtriebskraft:
\[
  F_{R,H,\max} = F_H.
\]
\textbf{Schritt 1:} Beide Kräfte einsetzen ($F_{R,H,\max}=\mu_H\cdot F_N$
und $F_H=F_G\cdot\sin(\alpha)$, siehe Kräfte Teil 1):
\[
  \mu_H\cdot F_N = F_G\cdot\sin(\alpha).
\]
\textbf{Schritt 2:} Normalkraft auf der schiefen Ebene einsetzen
($F_N=F_G\cdot\cos(\alpha)$, siehe Theoriebox zur Normalkraft):
\[
  \mu_H\cdot F_G\cdot\cos(\alpha) = F_G\cdot\sin(\alpha).
\]
\textbf{Schritt 3:} $F_G$ kürzen (auf beiden Seiten durch $F_G$ teilen):
\[
  \mu_H\cdot\cos(\alpha) = \sin(\alpha).
\]
\textbf{Schritt 4:} Durch $\cos(\alpha)$ teilen:
\[
  \boxed{\mu_H = \frac{\sin(\alpha)}{\cos(\alpha)} = \tan(\alpha)}.
\]
Am Grenzwinkel lässt sich der Haftreibungskoeffizient also allein aus dem
gemessenen Winkel bestimmen, ganz ohne Masse oder Gewichtskraft zu
kennen.
\end{theoriebox}

\begin{beispielbox}[Beispiel: Grenzwinkel messen]
Eine Holzkiste auf einer Holzunterlage beginnt bei einem Neigungswinkel
von $\alpha=\ang{27}$ zu rutschen.

\[
  \mu_H = \tan(\ang{27}) \approx 0.51.
\]
Das passt gut zum Tabellenwert für Holz auf Holz ($\mu_H\approx0.5$) aus
der Theoriebox zur Reibung.
\end{beispielbox}

\begin{questions}
\begin{taskbox}[Übung: Grenzwinkel auswerten]
\question[4] Eine Turnschuhsohle auf der Bodenmatte in der Turnhalle
beginnt erst bei einem Neigungswinkel von $\alpha=\ang{15}$ zu rutschen.

\begin{parts}
  \part[4] Berechne den Haftreibungskoeffizienten $\mu_H$ dieser
  Materialpaarung.
  \Antwortraum{4}
  \begin{solution}
  \[
    \mu_H = \tan(\ang{15}) \approx 0.27.
  \]
  \end{solution}
\end{parts}
\end{taskbox}
\end{questions}

% --- phyphox-Gruppenarbeit (M5/M6), wie in lernziele.md vorgeschlagen (L5).
% LZ9. -----------------------------------------------------------------------
\begin{groupbox}[phyphox-Experiment: Grenzwinkel und Reibungskoeffizienten messen]
Bearbeitet die folgende Aufgabe zu zweit oder in einer Kleingruppe.
Material: ein Buch oder Laptop als schiefe Ebene, ein Smartphone mit der
App phyphox (Sensor \glqq Inclination\grqq{} / Neigung), ein Testobjekt
(z.\,B. ein Radiergummi oder das Smartphone selbst, einmal mit und einmal
ohne Hülle).

\begin{parts}
  \part[2] \textbf{Aufbau:} Legt das Testobjekt auf die zunächst flache
  Unterlage (Buch/Laptop) und startet phyphox mit dem Sensor
  \glqq Inclination\grqq{} auf der Unterlage.
  \Antwortraum{2}

  \part[3] \textbf{Grenzwinkel messen:} Hebt die Unterlage sehr langsam
  an, bis das Testobjekt gerade zu rutschen beginnt, und lest in diesem
  Moment den Winkel $\alpha_H$ ab. Wiederholt dies für zwei verschiedene
  Bedingungen (z.\,B. mit und ohne Handyhülle) und berechnet jeweils
  $\mu_H=\tan(\alpha_H)$.
  \begin{center}
  \begin{tabular}{@{}lp{2.5cm}p{2.5cm}@{}}
  \toprule
  \textbf{Bedingung} & \textbf{$\alpha_H$} & \textbf{$\mu_H=\tan(\alpha_H)$} \\
  \midrule
  1 & & \\[0.6cm]
  2 & & \\[0.6cm]
  \bottomrule
  \end{tabular}
  \end{center}
  \Antwortraum{1}

  \part[3] \textbf{Gleitreibung messen:} Verkleinert nach Rutschbeginn
  den Winkel langsam wieder, bis sich das Testobjekt gerade noch mit
  konstanter Geschwindigkeit (nicht beschleunigt) weiterbewegt, lest
  diesen Winkel $\alpha_G$ ab und berechnet $\mu_G=\tan(\alpha_G)$ für
  beide Bedingungen.
  \begin{center}
  \begin{tabular}{@{}lp{2.5cm}p{2.5cm}@{}}
  \toprule
  \textbf{Bedingung} & \textbf{$\alpha_G$} & \textbf{$\mu_G=\tan(\alpha_G)$} \\
  \midrule
  1 & & \\[0.6cm]
  2 & & \\[0.6cm]
  \bottomrule
  \end{tabular}
  \end{center}
  \Antwortraum{1}

  \part[2] \textbf{Vergleich:} Vergleicht eure gemessenen Werte mit einem
  Literaturwert für eine ähnliche Materialpaarung (Internet- oder
  KI-Recherche, mit Quellenangabe). Stimmen eure Messungen ungefähr
  überein? Nennt eine mögliche Fehlerquelle bei diesem Experiment.
  \Antwortraum{2}
  \begin{solution}
  Individuelle Messwerte, von der Lehrperson auf $\mu_H>\mu_G$ und
  plausible Grössenordnung (z.\,B. Gummi/Kunststoff auf Buchdeckel eher
  im Bereich $\mu\approx0.3$ bis $0.6$) zu kontrollieren. Mögliche
  Fehlerquellen: ungleichmässiges Anheben, Ablesegenauigkeit des
  Neigungssensors, unterschiedlich raue Stellen der Unterlage.
  \end{solution}
\end{parts}
\end{groupbox}

% --- Abschluss: Clicker-Review über alle vier Kraftarten, wie in
% lernziele.md vorgeschlagen (Abschluss von L5). -----------------------------
\begin{questions}
\begin{taskbox}[Clicker-Review: alle vier Kraftarten]
\question[8] Ordne jede Aussage der passenden Kraftart zu: \textbf{(A)}
Gewichtskraft, \textbf{(B)} Federkraft, \textbf{(C)} Normalkraft,
\textbf{(D)} Reibungskraft.

\begin{parts}
  \part[2] \glqq Diese Kraft ist ortsabhängig: Auf dem Mond ist sie viel
  kleiner als auf der Erde, obwohl sich die Masse nicht ändert.\grqq{}
  \Antwortraum{2}
  \begin{solution}
  \textbf{A} (Gewichtskraft, $F_G=m\cdot g$).
  \end{solution}

  \part[2] \glqq Diese Kraft ist proportional zur Auslenkung aus der
  Ruhelage und wirkt immer rückstellend.\grqq{}
  \Antwortraum{2}
  \begin{solution}
  \textbf{B} (Federkraft, $F_H=-k\cdot x$, Hookesches Gesetz).
  \end{solution}

  \part[2] \glqq Diese Kraft ist eine Reaktionskraft der Unterlage und
  steht immer senkrecht zu deren Oberfläche.\grqq{}
  \Antwortraum{2}
  \begin{solution}
  \textbf{C} (Normalkraft, $F_N$, aus dem Kräftegleichgewicht senkrecht
  zur Unterlage).
  \end{solution}

  \part[2] \glqq Diese Kraft wirkt entgegen der Bewegungsrichtung (oder
  Bewegungstendenz) und ist proportional zur Normalkraft.\grqq{}
  \Antwortraum{2}
  \begin{solution}
  \textbf{D} (Reibungskraft, $F_R=\mu\cdot F_N$).
  \end{solution}
\end{parts}
\end{taskbox}
\end{questions}

\end{document}
